\subsection{Valores Mágicos}

\subsubsection{Regla: enumeradores y constantes en lugar de literales}

\textbf{Regla:} Prohibido usar valores numéricos o cadenas literales directamente en la lógica del juego; todo valor con significado en el dominio se gestiona mediante \texttt{enum} o constante de clase (Martin, 2008).

\textbf{Código correcto:}
\textit{Archivo \texttt{TurnState.cs}:}
\begin{lstlisting}[style=csharpstyle]
public enum TurnState
{
    WaitingForRoll,
    Moving,
    Finished
}
\end{lstlisting}

\textit{Archivo \texttt{TurnManager.cs}:}
\begin{lstlisting}[style=csharpstyle]
public sealed class TurnManager
{
    private const int MaximumPlayers = 4;

    public bool CanAddPlayer(int currentPlayerCount)
    {
        return currentPlayerCount < MaximumPlayers;
    }
}
\end{lstlisting}

\textbf{Código incorrecto:}
\begin{lstlisting}[style=csharpstyle]
public sealed class TurnManager
{
    public bool CanAddPlayer(int currentPlayerCount)
    {
        return currentPlayerCount < 4;
    }

    public void SetState(int state)
    {
        if (state == 2)
        {
            FinishTurn();
        }
    }

    private void FinishTurn()
    {
    }
}
\end{lstlisting}

\textbf{Aclaraciones:}
\begin{itemize}
    \item Los literales \texttt{0} y \texttt{1} se permiten como valores neutros: un contador que empieza en cero, un incremento o el primer índice. McConnell (2004) los señala como los únicos literales aceptables dentro de la lógica. Si el \texttt{0} o el \texttt{1} tiene significado en el dominio, como el primer turno de la partida, lleva nombre.
    \item Un literal se permite cuando es el valor de una definición con nombre: el inicializador de una constante, de una propiedad estática de solo lectura (Sección 2.3.2) o de un miembro de un \texttt{enum}. Esa definición es el lugar donde el valor recibe su nombre.
    \item Un valor que ya tiene nombre no se repite como literal en otro lugar, tampoco dentro de una cadena o de una expresión regular. La cadena se construye a partir de la constante, para que ambas no puedan quedar distintas.
\end{itemize}

\textbf{Código correcto (expresión regular):}
\begin{lstlisting}[style=csharpstyle]
public static class PlayerNameRules
{
    public const int MinimumLength = 3;
    public const int MaximumLength = 16;

    private static readonly Regex _namePattern = new Regex(
        $"^[A-Za-z0-9_]{{{MinimumLength},{MaximumLength}}}$");
}
\end{lstlisting}

\textbf{Código incorrecto (expresión regular):}
\begin{lstlisting}[style=csharpstyle]
public static class PlayerNameRules
{
    public const int MinimumLength = 3;
    public const int MaximumLength = 16;

    private static readonly Regex _namePattern = new Regex(@"^[A-Za-z0-9_]{3,16}$");
}
\end{lstlisting}
